%% kapa: exam=1 round=29 year=2018 period=1 form=A subject=economics subject_ko=경제학원론 seq=8 number=48
%% source: Q-Net 감정평가사 기출문제 (cst00309) artlSeq=5202808 file=extracted/1008/5202808/2018년 제29회 감정평가사 1차시험 기출문제/2018년 제29회 감정평가사 1차시험 기출문제/1교시 A형.pdf
%% conversion: kordoc 4.24.51 -> tools/gen1.py (auto); review=proofread-sonnet; flags=none
\begin{question}{48}
甲국 정부는 독점기업 A로 하여금 이윤극대화보다는 완전경쟁시장에서와 같이 사회적으로 효율적인 수준에서 생산하도록 규제하려고 한다. 사회적으로 효율적인 생산량이 달성되는 조건은? (단, 수요곡선은 우하향, 기업의 한계비용곡선은 우상향한다.)
\begin{choices}
  \item 평균수입 = 한계비용
  \item 평균수입 = 한계수입
  \item 평균수입 = 평균생산
  \item 한계수입 = 한계비용
  \item 한계수입 = 평균생산
\end{choices}
\end{question}
\answer{1}
